\documentclass[11pt,a4paper]{ctexart}
\usepackage[margin=2.2cm]{geometry}
\usepackage{xcolor}
\usepackage{enumitem}
\usepackage{amssymb}
\usepackage{booktabs}
\usepackage[normalem]{ulem}
\usepackage{array}
\usepackage{longtable}
\usepackage{titlesec}
\usepackage{hyperref}
\definecolor{navy}{HTML}{1A365D}
\hypersetup{colorlinks=true,linkcolor=navy,urlcolor=navy}
\setlist{nosep}
\titleformat{\section}{\Large\bfseries\color{navy}}{}{0pt}{}
\titleformat{\subsection}{\large\bfseries\color{navy}}{}{0pt}{}
\pagestyle{plain}
\renewcommand{\arraystretch}{1.25}

\title{\textbf{八年级上册语法总结}\\[4pt]
\large 人教版（新课纲版）八年级上册 · 按教材书末 Grammar 附录整理}
\author{初中英语资料站}
\date{2026-11-12 · 附配套专项练习 192 题}
\begin{document}
\maketitle
\section*{使用说明}
\begin{itemize}[leftmargin=1.4em]
\item 本讲义按八年级上册教材书末 Grammar 附录整理，表格与例句取自教材原文。
\item 网页版：\url{https://1901151458-prog.github.io/english-materials/g8a-grammar/}
\end{itemize}

\section{不定代词（Indefinite Pronouns）}
\noindent\textcolor{gray}{some / any、many / much、few / a few、复合不定代词}\par\smallskip
不指明代替任何特定名词或形容词的代词叫作不定代词。常见的有 \textbf{some、any、somebody、something、nobody、nothing、everybody、everything、one、another、other(s)、all、both、each、either、neither、none、many、much、few、a few、little、a little} 等。

不定代词大都可以代替名词和形容词，在句子中用作主语、宾语、定语等。需要注意以下几点：

\subsection*{1. some 和 any}

这两个代词既可代替或修饰可数名词，也可代替或修饰不可数名词。\textbf{some 多用于肯定句，any 则多用于否定句、疑问句和条件从句}。

\begin{itemize}[leftmargin=1.4em]
\item A: Are there \underline{any} apples in the fridge?\\ B: Yes, there are \underline{some}. / No, there aren’t \underline{any}.\\ 可数名词：any 用于疑问句，some / any 用于答句。
\item A: Is there \underline{any} water in the bottle?\\ B: Yes, there is \underline{some} water. / No, there isn’t \underline{any} water.\\ 不可数名词同样适用。
\end{itemize}

\begin{center}\fbox{\parbox{0.92\textwidth}{\small \textbf{教材提示：}在「Would you like some tea?」这类疑问句中则用 \textbf{some}，而不用 any，因为说话人希望得到对方\textbf{肯定的答复}。}}\end{center}

\subsection*{2. many 和 much}

这两个词均表示「很多」，但 \textbf{many 只代替或修饰可数名词复数}，而 \textbf{much 只代替或修饰不可数名词}。

\begin{itemize}[leftmargin=1.4em]
\item In the future, \underline{many} people will live in outer space.\\ people 可数复数，用 many。
\item There isn’t \underline{much} water left in the bottle.\\ water 不可数，用 much。
\end{itemize}

\subsection*{3. few 与 a few；little 和 a little}

\textbf{few / a few} 代替或修饰\textbf{可数名词复数}；\textbf{little / a little} 代替或修饰\textbf{不可数名词}。few 和 little 的意思是「很少；几乎没有」，侧重\textbf{否定}；a few 和 a little 的意思是「一些」，侧重\textbf{肯定}。

{\footnotesize
\begin{longtable}{p{3.35cm}p{3.35cm}p{3.35cm}p{3.35cm}}
\hline
词 & 修饰 & 含义倾向 & 例句 \\
\hline
\endfirsthead
\hline
词 & 修饰 & 含义倾向 & 例句 \\
\hline
\endhead
\hline
\endfoot
few & 可数复数 & 否定（几乎没有） & \textbf{Few} people know what happened. \\
a few & 可数复数 & 肯定（有一些） & We stayed in the old town for \textbf{a few} days. \\
little & 不可数 & 否定（几乎没有） & There is now \textbf{little} hope of success. \\
a little & 不可数 & 肯定（有一些） & Don’t worry. There is \textbf{a little} time left. \\
\end{longtable}
}

\subsection*{4. 复合不定代词}

由 some、any、no、every 加上 \textbf{-body、-one、-thing} 构成的不定代词叫作复合不定代词。除了 \textbf{no one} 以外，其他复合不定代词都写成一个词。复合不定代词只相当于名词，在句中用作主语、宾语或表语，\textbf{作主语时都作单数看待}。

\begin{itemize}[leftmargin=1.4em]
\item \underline{Everyone} says Yulong Mountain is beautiful.\\ 作主语，谓语用单数 says。
\item \underline{Nobody} knows what the future will be like.\\ 作主语，谓语用单数 knows。
\end{itemize}

\begin{center}\fbox{\parbox{0.92\textwidth}{\small \textbf{定语后置：}复合不定代词如果有定语修饰，\textbf{定语须置于其后}。\\ Did you do \textbf{anything fun} on your holiday? / I did \textbf{nothing special}. / Yesterday, I met \textbf{someone interesting}.}}\end{center}

\section{比较级与最高级：构成}
\noindent\textcolor{gray}{规则变化（表1）+ 不规则变化（表2）}\par\smallskip
大多数形容词和副词有三个等级：\textbf{原级}（原形）、\textbf{比较级}（「较……／更……」）、\textbf{最高级}（「最……」）。

{\footnotesize
\begin{longtable}{p{3.35cm}p{3.35cm}p{3.35cm}p{3.35cm}}
\hline
构成方法 & 原级 & 比较级 & 最高级 \\
\hline
\endfirsthead
\hline
构成方法 & 原级 & 比较级 & 最高级 \\
\hline
\endhead
\hline
\endfoot
单音节词和部分双音节词：一般加 -er / -est & tall / short / fast / high & taller / shorter / faster / higher & tallest / shortest / fastest / highest \\
以字母 e 结尾：加 -r / -st & large / late & larger / later & largest / latest \\
重读闭音节词尾只有一个辅音字母：双写再加 -er / -est & big / fat / thin & bigger / fatter / thinner & biggest / fattest / thinnest \\
以「辅音字母 + y」结尾的双音节词：变 y 为 i 再加 -er / -est & easy / heavy / early & easier / heavier / earlier & easiest / heaviest / earliest \\
多音节词和部分双音节词：词前加 more / most & beautiful / interesting / popular & more beautiful / more interesting / more popular & most beautiful / most interesting / most popular \\
\end{longtable}
}

{\footnotesize
\begin{longtable}{p{4.47cm}p{4.47cm}p{4.47cm}}
\hline
原级 & 比较级 & 最高级 \\
\hline
\endfirsthead
\hline
原级 & 比较级 & 最高级 \\
\hline
\endhead
\hline
\endfoot
good / well & better & best \\
bad / badly & worse & worst \\
many / much & more & most \\
little & less & least \\
far & farther / further & farthest / furthest \\
\end{longtable}
}

\section{比较级与最高级：用法}
\noindent\textcolor{gray}{表3 + as ... as 句型}\par\smallskip
{\footnotesize
\begin{longtable}{p{4.47cm}p{4.47cm}p{4.47cm}}
\hline
等级 & 用法 & 例句 \\
\hline
\endfirsthead
\hline
等级 & 用法 & 例句 \\
\hline
\endhead
\hline
\endfoot
比较级 & 表示两者（人或事物）的比较。在「比较级 + than」句型中，当 than 前后所使用的动词相同时，通常用助动词代替后面的动词，该助动词也可以省略。 & I’m \textbf{more outgoing than} my sister.\\ You know more about yourself \textbf{than I (do)}.\\ He works much harder than me.（口语中一般用 me） \\
最高级 & 表示三者或三者以上的比较，其中一个在某一方面超过其他几个。后面可用带有 of、in、among 的短语说明比较范围。\textbf{形容词最高级前一般要加定冠词 the}，副词最高级前的 the 可以省去。 & The blue whale is \textbf{the largest and heaviest} animal in the world.\\ In my country, July is \textbf{the hottest} month of the year.\\ The cheetah runs \textbf{(the) fastest} among all land animals. \\
\end{longtable}
}

\begin{center}\fbox{\parbox{0.92\textwidth}{\small \textbf{教材注释：}在形容词和副词的比较级前，可以用 \textbf{much、a little} 等修饰。\\ I’m feeling \textbf{much better} today. / My brother is \textbf{a little taller} than me.}}\end{center}

\subsection*{「as + 形容词或副词原形 + as」句型}

\textbf{（1）}表示两者在某一方面\textbf{相同}时，用「as + 原级 + as」。

\begin{itemize}[leftmargin=1.4em]
\item I think science is \underline{as important as} maths.\\ 我认为科学和数学一样重要。
\item Tom runs \underline{as fast as} Jack.\\ 汤姆跑得和杰克一样快。
\end{itemize}

\textbf{（2）}表示一方在某一方面\textbf{不及}另一方时，用「not as / so + 原级 + as」。

\begin{itemize}[leftmargin=1.4em]
\item It is \underline{not as / so warm} today as yesterday.\\ 今天不如昨天暖和。
\item He did \underline{not come as / so early} as Wang Lin.\\ 他没有王林来得那么早。
\end{itemize}

\section{情态动词（Modal Verbs）：can / could}
\noindent\textcolor{gray}{能力、请求允许、给予许可、提出请求和建议}\par\smallskip
情态动词 \textbf{can} 和 \textbf{could} 都可用于表示\textbf{能力、请求允许和给予许可}，以及\textbf{提出请求和建议}。

{\footnotesize
\begin{longtable}{p{6.70cm}p{6.70cm}}
\hline
功能 & 例句 \\
\hline
\endfirsthead
\hline
功能 & 例句 \\
\hline
\endhead
\hline
\endfoot
能力 & \textbf{Can} you speak English? 你会说英语吗？\\ He \textbf{could} sing very well when he was young. 他年轻时歌唱得很好。 \\
请求允许 & A: \textbf{Can} I put some flowers on the table? 我能在桌上摆些花吗？\\ \textbf{Could} I use your mobile phone? 我可以用你的手机吗？ \\
给予许可 & B: Yes, of course you \textbf{can}. 可以，当然可以。 \\
提出请求 & \textbf{Can} you help me hang up the photos? 你能帮我把照片挂起来吗？\\ \textbf{Could} you please sweep the floor? 请你打扫一下地板，好吗？ \\
提出建议 & We \textbf{can} eat in a restaurant this evening. 今晚我们可以在餐馆吃饭。\\ Maybe I \textbf{could} give you a lift to the supermarket. 也许我可以让你搭车去超市。\\ We \textbf{could} go fishing this weekend. 这个周末我们可以去钓鱼。 \\
\end{longtable}
}

\begin{center}\fbox{\parbox{0.92\textwidth}{\small \textbf{could 的特点：}表示请求时，could 比 can \textbf{更客气、更委婉}；但表示能力时，could 是 can 的过去式（He could sing very well when he was young）。}}\end{center}

\section{一般将来时（Simple Future Tense）}
\noindent\textcolor{gray}{will / shall + 动词原形；be going to}\par\smallskip
一般将来时通常由「\textbf{助动词 + 动词原形}」构成，表示将来某个时间要发生的动作或存在的状态，也表示将来经常或反复发生的动作，常与表示将来的时间状语连用，如 \textbf{tomorrow、next week、next year、in the future} 等。

\begin{itemize}[leftmargin=1.4em]
\item We \underline{will come} to see you next week.\\ 我们下周会来看你。
\item You \underline{will see} robots working everywhere in the future.\\ 将来你会看到机器人到处工作。
\item What \underline{will your life be like} in 20 years?\\ 20 年后你的生活会是什么样？
\item Some jobs today \underline{will disappear} in the coming years. However, there \underline{will be} more challenging jobs in the future.\\ 如今一些工作在未来几年会消失，但将来会有更有挑战性的工作。
\end{itemize}

\paragraph{will / shall}

\begin{itemize}[leftmargin=1.4em]
\item 口语中 will 常缩略为 \textbf{’ll}，will not 缩略为 \textbf{won’t}。
\item 表示「带意愿色彩的将来」时常用 will。
\item 主语为第一人称（I / we）时，也可用 \textbf{shall}（书面语更正式）。
\end{itemize}

\paragraph{be going to + 动词原形}

\begin{itemize}[leftmargin=1.4em]
\item 表示将要发生的动作或安排。
\item 也表示打算、计划、决定要做的事。
\item I’m \textbf{going to make} a sandwich.
\end{itemize}

\begin{center}\fbox{\parbox{0.92\textwidth}{\small \textbf{教材例句：}Shall we go to see the sci-fi film? / I shall write you a letter next month. / What are you going to do tonight?}}\end{center}

\section{句子种类（Sentence Types）}
\noindent\textcolor{gray}{陈述句 / 疑问句 / 祈使句 / 感叹句}\par\smallskip
英语句子按照用途可分为\textbf{陈述句、疑问句、祈使句和感叹句}四类。

{\footnotesize
\begin{longtable}{p{4.47cm}p{4.47cm}p{4.47cm}}
\hline
句子种类 & 用途 & 例句 \\
\hline
\endfirsthead
\hline
句子种类 & 用途 & 例句 \\
\hline
\endhead
\hline
\endfoot
陈述句 & 用于说明事实或说话人的看法。 & Bamboo is a symbol of goodness in Chinese culture.\\ I think you are going to be a great writer. \\
疑问句 & 用于提出问题。 & Did you do anything special on holiday?\\ What do you want to be when you grow up?\\ Which would you like, tea or coffee?\\ It’s a lovely day, isn’t it? \\
祈使句 & 用于表示请求、命令、劝说、号召、警告、建议等。 & Let’s make some mashed potatoes for lunch.\\ Don’t be nervous.\\ Pay attention to your body language.\\ Come in, please. \\
感叹句 & 用于表达强烈的感情。 & How clever you are!\\ What wonderful memories I have! \\
\end{longtable}
}

\begin{center}\fbox{\parbox{0.92\textwidth}{\small \textbf{感叹句两种开头：}\textbf{How + 形容词/副词}（How clever you are!）／\textbf{What + 名词短语}（What wonderful memories I have!）。}}\end{center}

\section{句子的成分与五种基本句型}
\noindent\textcolor{gray}{七大成分（表5）+ 五种句型（表6）}\par\smallskip
句子的成分可分为\textbf{主语、谓语、表语、宾语、宾语补足语、定语和状语}等。主语和谓语是句子的主体部分。

{\footnotesize
\begin{longtable}{p{4.47cm}p{4.47cm}p{4.47cm}}
\hline
句子成分 & 意义 & 例句 \\
\hline
\endfirsthead
\hline
句子成分 & 意义 & 例句 \\
\hline
\endhead
\hline
\endfoot
主语 & 表示句子所说的是「什么人」或「什么事物」，一般由名词、代词或相当于名词的词或短语等充当。 & My sister wants to be a singer.\\ Everything tasted good.\\ The film starts at half past two. \\
谓语 & 说明主语「做什么」「是什么」或者「怎么样」。谓语用动词，和主语在人称和数两方面必须一致。 & I love our new house.\\ I often do my homework after school.\\ Many people are singing and dancing in the park. \\
表语 & 说明主语「是什么」或者「怎么样」，由名词、形容词或相当于名词或形容词的词或短语等充当，和系动词一起构成谓语。 & The metro stations are beautiful.\\ You seem positive about the future.\\ A good friend is like a mirror. \\
宾语 & 表示动作、行为的对象，由名词、代词或相当于名词的词或短语等充当。 & We need some milk and butter.\\ Could you take me to the flower shop? \\
宾语补足语 & 用来补足宾语的意义，一般置于宾语之后，由名词、形容词、不定式、分词、介词短语等充当。 & I found the book very interesting.\\ My parents want me to be a doctor. \\
定语 & 用来修饰名词或代词。作定语的除形容词外，还有代词、数词、名词、介词短语等。 & What’s your name?\\ I am going to make three resolutions.\\ Bees are hard workers, and their greatest work is pollination. \\
状语 & 用来修饰动词、形容词或副词，一般表示行为发生的时间、地点、目的、方式、程度等。 & I like it very much.\\ We will meet at nine o’clock.\\ My friend will go to the airport to pick you up. \\
\end{longtable}
}

{\footnotesize
\begin{longtable}{p{3.35cm}p{3.35cm}p{3.35cm}p{3.35cm}}
\hline
类别 & 主语 S & 谓语部分 & 例句 \\
\hline
\endfirsthead
\hline
类别 & 主语 S & 谓语部分 & 例句 \\
\hline
\endhead
\hline
\endfoot
主谓 (SV) & Ella & is dancing.（不及物动词） & Ella is dancing. \\
主（系）表 (SP) & They & are / look（系动词）+ 表语 & They are friends. / They look the same. \\
主谓宾 (SVO) & I & like（及物动词）+ 宾语 & I like bananas. \\
主谓宾宾 (SVOO) & She & gave（及物动词）+ 间接宾语 me + 直接宾语 a pen & She gave me a pen. \\
主谓宾补 (SVOC) & He & made（及物动词）+ 宾语 us + 宾补 laugh & He made us laugh. \\
\end{longtable}
}

\begin{center}\fbox{\parbox{0.92\textwidth}{\small \textbf{记忆法：}五种句型 = 「谁 + 做什么」的五种扩展：\textbf{SV}（只有动作）→ \textbf{SP}（是什么样）→ \textbf{SVO}（做什么）→ \textbf{SVOO}（给谁什么）→ \textbf{SVOC}（让谁怎么样）。}}\end{center}

\section{配套专项练习}
\noindent\textcolor{gray}{四个专项各 48 题、共 192 题，支持在线作答与即时解析}\par\smallskip
\paragraph{不定代词专项 · 48 题}

\paragraph{比较级与最高级专项 · 48 题}

\paragraph{一般将来时与情态动词专项 · 48 题}

\paragraph{句子种类与句子成分专项 · 48 题}

\begin{center}\fbox{\parbox{0.92\textwidth}{\small \textbf{建议顺序：}先按本页 01–07 讲通规则 → 再进入对应专项练习自测 → 错题回到小节复习。四个专项合计 \textbf{192 题}。\\ \textbf{考前速查：}一页纸速查表 →（A4 单页，可打印）}}\end{center}

\end{document}
